\section{Notation}

We use $U$ to denote an open set in $\R^n$, and $M$ to denote a compact Riemannian manifold of dimension $n \ge 3$ with metric $g$. We use $\dd{V}_g$ to denote the volume form on $M$ induced by $g$. We use $\Delta$ to denote the Laplacian, and $\nabla$ to denote the covariant derivative. We use $C$ to denote a constant that may change from line to line, but does not depend on any variables.

\subsection{Linear Algebra}

\begin{enumerate}
  \item We use $\langle\cdot, \cdot\rangle$ to denote the inner product on $\R^n$ and the inner product on $T_xM$ induced by $g$.
  \item We write $A = ((a_{ij}))$ to mean $A$ is an $m \times n$ matrix with $(i, j)^{\text{th}}$ entry $a_{ij}$.
  \item If $A = ((a_{ij}))$ and $B = ((b_{ij}))$ are $m \times n$ matrices, then:
    \[%
      A : B = \sum_{i=1}^m \sum_{j=1}^n a_{ij} b_{ij}
    ,\]%
    and
    \[%
      |A| = (A : A)^{1/2} = \left(\sum_{i=1}^m \sum_{j=1}^n a_{ij}^2\right)^{1/2}
    .\]%
  \item We denote the space of invertible real matrices by $\GL(n, \R)$ (or $\GL(n, \C)$ for complex matrices).
  \item We denote the space of real symmetric matrices by $\S(n, \R)$ (or $\S(n, \C)$ for complex symmetric matrices).
\end{enumerate}

\subsection{Functions}

\begin{enumerate}
  \item If $u : U \to \R$, we write $u(x) = u(x_1, \cdots, x_n)$.
  \item If $\u : U \to \R^m$, we write $\u(x) = (u^1(x), \cdots, u^m(x))$.
  \item If $\Sigma$ is a smooth $(n-1)$-dimensional surface in $\R^m$, we write:
    \[%
      \int_\Sigma f \dd{S}
    ,\]%
    for the integral of $f$ over $\Sigma$, with respect to the $(n-1)$-dimensional surface measure. If $C$ is a curve in $\R^m$, we denote it by:
    \[%
      \int_C f \dd{l}
    ,\]%
    the integral of $f$ over $C$ with respect to arclength.
  \item Averages:
    \[%
      \dashint_{B(x,r)} f \dy = \frac{1}{\alpha(n)r^n} \int_{B(x,r)} f \dy =~\text{average of $f$ over the ball $B(x,r)$}
    ,\]%
    and
    \[%
      \dashint_{\partial B(x,r)} f \dd{S} = \frac{1}{n\alpha(n)r^{n-1}} \int_{\partial B(x,r)} f \dd{S} =~\text{average of $f$ over the sphere $\partial B(x,r)$}
    .\]%
\end{enumerate}

\subsection{Derivatives}

\begin{enumerate}
  \item I write:
    \[%
      \pdv{u}{x_i} = \pd{x_i} u = u_{x_i}
    .\]%
  \item Similarly, we also have:
    \[%
      \pdv{u}{x_i,x_j} = \pd[style-var=single]{x_i,x_j} u = u_{x_i,x_j}
    .\]%
  \item Given a multi-index $\alpha$, define:
    \[%
      D^\alpha u(x) = \pdv[order={\alpha_1, \alpha_n}, mixed-order={|\alpha|}, sep-ord-inf={\cdots}]{u(x)}{x_1, x_n} = \pd[\alpha_1]{x_1} \cdots \pd[\alpha_n]{x_n} u(x)
    .\]%
    This is the standard multi-index notation for partial derivatives.
  \item If we specify a specific variable, we write:
    \[%
      D_xu = (u_{x_1}, \cdots, u_{x_n}), \quad D_yu = (u_{y_1}, \cdots, u_{y_n})
    .\]%
  \item If $k$ is a positive integer, the set of all partial derivatives of order $k$ is denoted by:
    \[%
      D^ku(x) := \{D^\alpha u(x) : |\alpha| = k\}
    .\]%
\end{enumerate}

